@ID: OW20D1P1I
@BIDANG: Analisis Real

Jika himpunan $B\subseteq\mathbb R$ tidak kosong dan terbatas di bawah, serta himpunan

$$
A=\{x\mid x \text{ adalah batas bawah dari } B\},
$$

maka $\sup A=\ldots$

@ID: OW20D1P2I
@BIDANG: Analisis Real

Diketahui fungsi $f$ terintegral pada $[a,b]$, dengan fungsi primitif $F:[a,b]\to\mathbb R$ yang didefinisikan oleh

$$
F(x)=\int_a^x f(t)\,dt
$$

untuk setiap $x\in[a,b]$.

Jika

$$
E=\{x\in[a,b]:F'(x)\ne f(x)\},
$$

maka closure dari $[a,b]\setminus E$ sama dengan $\ldots$

@ID: OW20D1P3I
@BIDANG: Struktur Aljabar

Misalkan $A$ merupakan suatu gelanggang polinom dengan koefisien di bilangan kompleks $\mathbb C$. Penjumlahan dua unsur di $A$ sama dengan penjumlahan pada gelanggang polinom biasa. Kemudian, perkalian $*$ di $A$ didefinisikan oleh

$$
(\beta x^k)*(\alpha x^j)=\beta\,\overline{\alpha}\,x^{k+j},
$$

untuk setiap $\alpha,\beta\in\mathbb C$.

Hasil dari

$$
(x^{2020}+2ix^7+9)*((i-1)x^{20}+(1-i))
$$

adalah $\ldots$

@ID: OW20D1P4I
@BIDANG: Struktur Aljabar

Misalkan $G$ adalah suatu grup yang dibangun oleh dua unsur $a$ dan $b$. Salah satu relasi yang dipenuhi oleh $a$ dan $b$ adalah

$$
a^2=b^2=(ab)^4=1.
$$

Maka orde terbesar dari $G$ yang mungkin adalah $\ldots$

@ID: OW20D1P5I
@BIDANG: Struktur Aljabar

Misalkan

$$
B=
\begin{pmatrix}
\mathbb Z_{2020} & 0\\
0 & \mathbb Z_{2020}
\end{pmatrix}
=
\left\{
\begin{pmatrix}
a&0\\
0&b
\end{pmatrix}
:a,b\in\mathbb Z_{2020}
\right\}
$$

merupakan suatu gelanggang. Jika

$$
M=\{\alpha\in B:\exists n\in\mathbb N\text{ sehingga }\alpha^n=0\},
$$

maka $M=\ldots$

@ID: OW20D1P6I
@BIDANG: Kombinatorika

Banyaknya solusi bilangan bulat tak negatif dari persamaan

$$
a+b+c=13
$$

dengan syarat $a\le4$, $2\le b\le5$, dan $1\le c\le6$ adalah $\ldots$

@ID: OW20D2P1I
@BIDANG: Aljabar Linear

Diketahui matriks

$$
A=
\begin{pmatrix}
1&0&2\\
-1&a&10\\
0&1&b
\end{pmatrix}
$$

memiliki rank sama dengan $2$. Nilai $ab$ adalah $\ldots$

@ID: OW20D2P2I
@BIDANG: Aljabar Linear

Misalkan $A$ sebuah matriks berukuran $9\times9$ yang memenuhi sifat:

(a) Semua komponen baris pertama matriks $A$ berbeda, yaitu bilangan asli $\{n_1,n_2,\ldots,n_9\}$.

(b) Komponen setiap baris lainnya merupakan suatu permutasi dari komponen pada baris pertama.

Nilai eigen yang senantiasa dimiliki oleh matriks $A$ adalah $\ldots$

@ID: OW20D2P3I
@BIDANG: Aljabar Linear

Solusi dari relasi rekuren

$$
u_n=3u_{n-2}-2u_{n-3},\qquad n\ge3,
$$

dengan syarat awal $u_0=1$, $u_1=0$, dan $u_2=0$ adalah $\ldots$

@ID: OW20D2P4I
@BIDANG: Analisis Kompleks

Banyak bilangan kompleks $z$ sehingga

$$
z^{2020}=1
$$

tetapi $z^{20}\ne1$ adalah $\ldots$

@ID: OW20D2P5I
@BIDANG: Analisis Kompleks

Cakram terbuka

$$
A=\left\{z\in\mathbb C:\left|z-\frac13\right|<r\right\}
$$

dipetakan oleh fungsi

$$
f(z)=\frac{z}{z+1}
$$

menjadi cakram terbuka

$$
B=\left\{z\in\mathbb C:|z|<\frac12\right\}.
$$

Nilai $r$ adalah $\ldots$

@ID: OW20D2P6I
@BIDANG: Analisis Kompleks

Penyelesaian dari persamaan

$$
ie^z+1=0
$$

yang memenuhi $4<|z|<5$ adalah $\ldots$

@ID: OW20D1P1U
@BIDANG: Analisis Real

Diberikan fungsi $f:[-1,1]\to\mathbb R$ dengan

$$
f(x)=
\begin{cases}
0,&-1\le x\le0,\\
1,&0<x\le1.
\end{cases}
$$

Selanjutnya diberikan fungsi $g:[-1,1]\to\mathbb R$ yang memenuhi $g'(x)=f(x)$ untuk setiap $x\in[-1,1]\setminus\{0\}$.

Apakah ada fungsi kontinu $\varphi:[-1,1]\to\mathbb R$ sehingga

$$
g(x)=\varphi(x)(x+1)
$$

untuk setiap $x\in[-1,1]$? Jika ada, tentukan fungsi $\varphi$ tersebut. Berikan penjelasan pada jawaban Saudara.

@ID: OW20D1P2U
@BIDANG: Analisis Real

Diketahui $x_1=\frac12$ dan

$$
x_{n+1}-x_n=\frac{n+1}{n+2},
$$

untuk setiap $n\in\mathbb N$.

Selidiki apakah barisan $(x_n)$ konvergen. Jika konvergen, tentukan limitnya. Berikan penjelasan pada jawaban Saudara.

@ID: OW20D1P3U
@BIDANG: Struktur Aljabar

Misalkan $G$ adalah suatu grup hingga yang tidak mengandung unsur berorde $3$. Misalkan untuk setiap $x,y\in G$ berlaku

$$
(xy)^3=x^3y^3.
$$

(a) Buktikan bahwa untuk setiap $g\in G$ terdapat $z\in G$ sehingga $g=z^3$.

(b) Tunjukkan bahwa untuk setiap $x,y\in G$ berlaku $xy^2=y^2x$.

(c) Buktikan bahwa $G$ komutatif.

@ID: OW20D1P4U
@BIDANG: Kombinatorika

Diberikan himpunan $A=\{a_1,a_2,\ldots,a_n\}$. Bila Anda memilih $2^{n-1}+1$ himpunan bagian berbeda dari $A$, buktikan bahwa di antara himpunan-himpunan bagian terpilih tersebut terdapat dua himpunan sehingga yang satu merupakan himpunan bagian dari yang lain.

@ID: OW20D2P1U
@BIDANG: Aljabar Linear

Tentukan bilangan bulat positif $k$ terkecil sedemikian sehingga untuk setiap $k$ buah matriks $A_1,A_2,\ldots,A_k$ berukuran $100\times100$ dengan komponen real, senantiasa terdapat bilangan real $\alpha_1,\alpha_2,\ldots,\alpha_k$ yang tidak semuanya nol sehingga

$$
\det(\alpha_1A_1+\alpha_2A_2+\cdots+\alpha_kA_k)=0.
$$

Jelaskan jawaban Anda.

@ID: OW20D2P2U
@BIDANG: Aljabar Linear

Misalkan

$$
e_1=\begin{pmatrix}1\\0\\0\end{pmatrix},\quad
e_2=\begin{pmatrix}0\\1\\0\end{pmatrix},\quad
e_3=\begin{pmatrix}0\\0\\1\end{pmatrix}.
$$

Misalkan juga $b$ sebuah vektor tak nol di $\mathbb R^3$, dan diketahui bahwa tiga vektor $e_1$, $2e_1+e_2$, dan $e_1+e_2+e_3$ merupakan solusi suatu sistem persamaan linear $Ax=b$, dengan $A$ suatu matriks berukuran $3\times3$.

(a) Buktikan bahwa kombinasi linear $a_1e_1+a_2e_2+a_3e_3$, dengan $a_1,a_2,a_3$ skalar, juga merupakan solusi sistem persamaan linear $Ax=b$ jika dan hanya jika $a_1+a_3=a_2$.

(b) Jika

$$
b=
\begin{pmatrix}
b_1\\
b_2\\
b_3
\end{pmatrix},
$$

tentukan semua nilai eigen real dari $A$.

@ID: OW20D2P3U
@BIDANG: Analisis Kompleks

Misalkan $z\in\mathbb C$ sehingga

$$
|z|+|z-2020|=2020.
$$

Tunjukkan bahwa

$$
|z-(20+20i)|\ge20.
$$

@ID: OW20D2P4U
@BIDANG: Kombinatorika

Untuk bilangan bulat $x,y\ge1$, tuliskan bukti kombinatorial dari

$$
xy=\binom{x+y}{2}-\binom{x}{2}-\binom{y}{2}.
$$